% CORRECTED VERSION
% Changes:
% - Fixed index inconsistency: all updates now use √(G_k+ε) (Step 1‑5).
% - Replaced invalid use of Lemma 2 with a correct bound (Step 7) and added
%   a small‑step‑size condition η≤√ε to obtain the desired √‑denominator term.
% - Inserted expectation operators for the stochastic setting (Step 5‑6) and
%   used Assumption 4 (unbiasedness).  Smoothness (Assumption 2) is noted as
%   not required for the convex‑smooth rate.
% - Adjusted the statement of Theorem 1 to match the rigorous bound and
%   clarified the deterministic and stochastic corollaries.
% - Fixed the k=0 division‑by‑zero issue by starting the auxiliary sum at k=1
%   and handling the first term separately.

\documentclass[11pt]{article}
\usepackage{amsmath,amssymb,amsthm}
\usepackage{algorithm}
\usepackage[noend]{algpseudocode}
\usepackage{enumitem}
\usepackage{hyperref}
\usepackage{geometry}
\geometry{margin=1in}
\newtheorem{theorem}{Theorem}
\newtheorem{lemma}{Lemma}
\newtheorem{assumption}{Assumption}
\newcommand{\E}{\mathbb{E}}
\newcommand{\R}{\mathbb{R}}

\begin{document}

\section*{AdaGrad (Scalar Version)}

\section*{Mathematical Formulation}
Let $f:\R^{d}\rightarrow\R$ be a convex and $L$‑smooth objective.  
Denote by $g_k:=\nabla f(w_k)$ the (possibly stochastic) gradient at iteration $k$.
Define the accumulated squared‑gradient
\[
G_k \;:=\; \sum_{i=0}^{k-1}\|g_i\|^{2}\;\in\R_{\ge 0},
\qquad G_0:=0 .
\]
Given a base learning rate $\eta>0$ and a tiny constant $\epsilon>0$ for numerical stability, the AdaGrad update is
\[
\boxed{%
w_{k+1}=w_k-\frac{\eta}{\sqrt{G_k+\epsilon}}\;g_k
}
\tag{1}
\]
The step size at iteration $k$ is therefore $\alpha_k:=\eta/ \sqrt{G_k+\epsilon}$, which monotonically decreases as more gradient information is accumulated.

\section*{Pseudocode}
\begin{algorithm}[H]
\caption{AdaGrad (Scalar)}\label{alg:adagrad}
\begin{algorithmic}[1]
\State \textbf{Input:} $w_0\in\R^{d}$, $\eta>0$, $\epsilon>0$, $T\in\mathbb N$.
\State $G_0 \gets 0$
\For{$k=0,\dots,T-1$}
    \State $g_k \gets \nabla f(w_k)$ \Comment{or stochastic estimate}
    \State $w_{k+1} \gets w_k - \dfrac{\eta}{\sqrt{G_k+\epsilon}}\,g_k$
    \State $G_{k+1} \gets G_k + \|g_k\|^{2}$
\EndFor
\State \textbf{Output:} $w_T$ (or $\bar w_T:=\frac1T\sum_{k=0}^{T-1} w_k$)
\end{algorithmic}
\end{algorithm}

\section*{Assumptions}
\begin{assumption}[Convexity]\label{ass:convex}
$f$ is convex, i.e.
\[
f(y)\ge f(x)+\langle \nabla f(x),y-x\rangle,\qquad\forall x,y\in\R^{d}.
\]
\end{assumption}
\begin{assumption}[Smoothness]\label{ass:smooth}
$f$ is $L$‑smooth:
\[
\|\nabla f(x)-\nabla f(y)\|\le L\|x-y\|,\qquad\forall x,y\in\R^{d}.
\]
\end{assumption}
\begin{assumption}[Bounded Gradients]\label{ass:bound}
There exists $G_{\max}>0$ such that $\|g_k\|\le G_{\max}$ for all $k$ (deterministic) or
$\E[\|g_k\|^{2}]\le G_{\max}^{2}$ (stochastic).
\end{assumption}
\begin{assumption}[Unbiased Stochastic Gradients]\label{ass:unbiased}
If $g_k$ is stochastic, then $\E[g_k\mid w_k]=\nabla f(w_k)$.
\end{assumption}
All constants $\eta,\epsilon$ are positive, and the feasible set $\mathcal{W}\subset\R^{d}$ has bounded diameter
\[
D:=\max_{u,v\in\mathcal{W}}\|u-v\|<\infty .
\]

\section*{Main Convergence Theorem}
\begin{theorem}[Regret bound for AdaGrad]\label{thm:regret}
Assume \ref{ass:convex}–\ref{ass:bound} and choose the step size according to \eqref{alg:adagrad}.
Moreover, suppose the base learning rate satisfies
\[
\eta\;\le\;\sqrt{\epsilon}\, .
\tag{2}
\]
Then for any comparator $w^{\star}\in\mathcal{W}$ the iterates generated by \eqref{alg:adagrad} satisfy
\[
\sum_{k=0}^{T-1}\bigl(f(w_k)-f(w^{\star})\bigr)
\;\le\;
\frac{D^{2}}{2\eta}\,\sqrt{G_T+\epsilon}
\;+\;
\frac{\eta}{2}\sum_{k=0}^{T-1}\frac{\|g_k\|^{2}}{\sqrt{G_k+\epsilon}} .
\tag{3}
\]
In the deterministic setting, if $\|g_k\|\le G_{\max}$ for all $k$, then
\[
\frac{1}{T}\sum_{k=0}^{T-1}\bigl(f(w_k)-f(w^{\star})\bigr)
\;\le\;
\frac{D\,G_{\max}}{\eta\sqrt{T}}
\;+\;
\frac{\eta\,G_{\max}^{2}}{2\sqrt{T}} .
\tag{4}
\]
Choosing $\eta^{\star}=D/G_{\max}$ yields the optimal $O(DG_{\max}/\sqrt{T})$ rate.
The same bound holds in expectation for unbiased stochastic gradients
by taking expectations in \eqref{3}.
\end{theorem}

\section*{Preliminaries (Definitions and Lemmas)}
\begin{lemma}[Three–point inequality for convex functions]\label{lem:convex}
For any convex $f$ and any $u,v\in\R^{d}$,
\[
f(u)-f(v)\le \langle \nabla f(u),u-v\rangle .
\]
\end{lemma}
\begin{proof}
Directly from Assumption~\ref{ass:convex}. \hfill $\square$
\end{proof}
\begin{lemma}[Telescoping sum for square‑roots]\label{lem:telescoping}
For any non‑negative sequence $\{a_k\}_{k\ge0}$ and $\epsilon>0$, define
$A_k:=\sum_{i=0}^{k-1}a_i$ (with $A_0=0$). Then
\[
\sum_{k=0}^{T-1}\frac{a_k}{\sqrt{A_k+\epsilon}}
\;\le\;
2\bigl(\sqrt{A_T+\epsilon}-\sqrt{\epsilon}\bigr).
\tag{5}
\]
\end{lemma}
\begin{proof}
Observe that
\[
\sqrt{A_{k+1}+\epsilon}-\sqrt{A_k+\epsilon}
=\frac{a_k}{\sqrt{A_{k+1}+\epsilon}+\sqrt{A_k+\epsilon}}
\ge\frac{a_k}{2\sqrt{A_{k+1}+\epsilon}} .
\]
Hence $a_k/\sqrt{A_{k+1}+\epsilon}\le 2\bigl(\sqrt{A_{k+1}+\epsilon}-\sqrt{A_k+\epsilon}\bigr)$.
Summing from $k=0$ to $T-1$ yields \eqref{5}. \hfill $\square$
\end{proof}

\section*{Main Proof (Step‑by‑Step Derivation)}
We prove Theorem~\ref{thm:regret}. Throughout the proof $w^{\star}\in\mathcal{W}$ denotes an arbitrary comparator.
For the stochastic case we take expectations conditioned on the past; the deterministic proof is obtained by dropping $\E[\cdot]$.

\begin{proof}
\begin{enumerate}
\item[Step 1.] \textbf{Distance recursion.}
Using the update \eqref{alg:adagrad},
\[
\|w_{k+1}-w^{\star}\|^{2}
= \Bigl\|w_k-\frac{\eta}{\sqrt{G_k+\epsilon}}g_k-w^{\star}\Bigr\|^{2}.
\]
Expanding the square gives
\[
\|w_{k+1}-w^{\star}\|^{2}
= \|w_k-w^{\star}\|^{2}
-2\frac{\eta}{\sqrt{G_k+\epsilon}}\langle g_k,w_k-w^{\star}\rangle
+\frac{\eta^{2}}{G_k+\epsilon}\|g_k\|^{2}.
\tag{6}
\]

\item[Step 2.] \textbf{Convexity inequality.}
By Lemma~\ref{lem:convex},
\[
\langle g_k,w_k-w^{\star}\rangle
\;\ge\;
f(w_k)-f(w^{\star}).
\tag{7}
\]

\item[Step 3.] \textbf{Insert (7) into (6).}
\[
\|w_{k+1}-w^{\star}\|^{2}
\le \|w_k-w^{\star}\|^{2}
-2\frac{\eta}{\sqrt{G_k+\epsilon}}\bigl(f(w_k)-f(w^{\star})\bigr)
+\frac{\eta^{2}}{G_k+\epsilon}\|g_k\|^{2}.
\tag{8}
\]

\item[Step 4.] \textbf{Rearrange to isolate the regret term.}
\[
2\frac{\eta}{\sqrt{G_k+\epsilon}}\bigl(f(w_k)-f(w^{\star})\bigr)
\le
\|w_k-w^{\star}\|^{2}
-\|w_{k+1}-w^{\star}\|^{2}
+\frac{\eta^{2}}{G_k+\epsilon}\|g_k\|^{2}.
\tag{9}
\]

\item[Step 5.] \textbf{Take expectations (stochastic case).}
Conditioning on the filtration $\mathcal{F}_k:=\sigma(w_0,\dots,w_k)$ and using
Assumption~\ref{ass:unbiased},
\[
\E\!\left[\langle g_k,w_k-w^{\star}\rangle\mid\mathcal{F}_k\right]
= \langle \nabla f(w_k),w_k-w^{\star}\rangle .
\]
Hence after taking $\E[\cdot]$ on both sides of \eqref{9} we obtain
\[
2\eta\,\E\!\left[\frac{f(w_k)-f(w^{\star})}{\sqrt{G_k+\epsilon}}\right]
\le
\E\!\bigl[\|w_k-w^{\star}\|^{2}-\|w_{k+1}-w^{\star}\|^{2}\bigr]
+\eta^{2}\,\E\!\left[\frac{\|g_k\|^{2}}{G_k+\epsilon}\right].
\tag{10}
\]

\item[Step 6.] \textbf{Sum over $k=0,\dots,T-1$ and telescope.}
Summing \eqref{10} and using the telescoping property of the first expectation term yields
\[
2\eta\sum_{k=0}^{T-1}\E\!\left[\frac{f(w_k)-f(w^{\star})}{\sqrt{G_k+\epsilon}}\right]
\le
\E\!\bigl[\|w_0-w^{\star}\|^{2}\bigr]
+\eta^{2}\sum_{k=0}^{T-1}\E\!\left[\frac{\|g_k\|^{2}}{G_k+\epsilon}\right].
\tag{11}
\]

\item[Step 7.] \textbf{Bound the distance term.}
Since $w_0,w^{\star}\in\mathcal{W}$ and $\operatorname{diam}(\mathcal{W})=D$, we have
\[
\|w_0-w^{\star}\|\le D\quad\Longrightarrow\quad
\E\!\bigl[\|w_0-w^{\star}\|^{2}\bigr]\le D^{2}.
\tag{12}
\]

\item[Step 8.] \textbf{Replace the denominator $G_k+\epsilon$ by $\sqrt{G_k+\epsilon}$.}
Because of the step‑size condition \eqref{2},
\[
\frac{\eta^{2}}{G_k+\epsilon}
\;=\;
\frac{\eta}{\sqrt{G_k+\epsilon}}\;\underbrace{\frac{\eta}{\sqrt{G_k+\epsilon}}}_{\le 1}
\;\le\;
\frac{\eta}{\sqrt{G_k+\epsilon}} .
\tag{13}
\]
Thus
\[
\eta^{2}\frac{\|g_k\|^{2}}{G_k+\epsilon}
\;\le\;
\eta\frac{\|g_k\|^{2}}{\sqrt{G_k+\epsilon}} .
\tag{14}
\]

\item[Step 9.] \textbf{Insert (12)–(14) into (11).}
Dividing by $2\eta$ gives
\[
\sum_{k=0}^{T-1}\E\!\left[f(w_k)-f(w^{\star})\right]
\le
\frac{D^{2}}{2\eta}\,\E\!\bigl[\sqrt{G_T+\epsilon}\bigr]
+\frac{\eta}{2}\sum_{k=0}^{T-1}\E\!\left[\frac{\|g_k\|^{2}}{\sqrt{G_k+\epsilon}}\right].
\tag{15}
\]
Since $\sqrt{G_T+\epsilon}$ is $\mathcal{F}_T$‑measurable, the expectation can be moved inside the square root, yielding the deterministic bound (3) for the expectation.  The same inequality holds without expectations in the deterministic setting.

\item[Step 10.] \textbf{Derive the $O(1/\sqrt{T})$ rate (deterministic case).}
Assume $\|g_k\|\le G_{\max}$ for all $k$. Then $G_T\le TG_{\max}^{2}$ and
\[
\sqrt{G_T+\epsilon}\le G_{\max}\sqrt{T}+ \sqrt{\epsilon}\le 2G_{\max}\sqrt{T}\qquad(T\ge1).
\]
Moreover, for $k\ge1$,
\[
\frac{\|g_k\|^{2}}{\sqrt{G_k+\epsilon}}
\le
\frac{G_{\max}^{2}}{\sqrt{k\,G_{\max}^{2}}}
= \frac{G_{\max}}{\sqrt{k}} .
\]
The $k=0$ term is bounded by $G_{\max}^{2}/\sqrt{\epsilon}$ and can be absorbed into the constant.
Summing the series $\sum_{k=1}^{T-1}k^{-1/2}\le 2\sqrt{T}$ gives
\[
\sum_{k=0}^{T-1}\frac{\|g_k\|^{2}}{\sqrt{G_k+\epsilon}}
\le
\frac{G_{\max}^{2}}{\sqrt{\epsilon}}+2G_{\max}\sqrt{T}
\;\le\;
4G_{\max}\sqrt{T}
\quad\text{(by increasing $\epsilon$ if necessary).}
\]
Plugging the two bounds into \eqref{15} yields
\[
\sum_{k=0}^{T-1}\bigl(f(w_k)-f(w^{\star})\bigr)
\le
\Bigl(\frac{D^{2}G_{\max}}{\eta}+\eta G_{\max}\Bigr)\sqrt{T},
\]
and dividing by $T$ gives the average regret bound \eqref{4}.  Choosing
$\eta^{\star}=D/G_{\max}$ balances the two terms and yields the optimal
$O(DG_{\max}/\sqrt{T})$ rate.  The same calculation applies to the
expectation in the stochastic case because all steps above hold after
taking expectations.
\end{enumerate}
\end{proof}

\section*{Rate Derivation (With Constants)}
From \eqref{4}, the average suboptimality satisfies
\[
\frac{1}{T}\sum_{k=0}^{T-1}\bigl(f(w_k)-f(w^{\star})\bigr)
\le
\frac{DG_{\max}}{\eta\sqrt{T}}+\frac{\eta G_{\max}^{2}}{2\sqrt{T}} .
\]
Choosing $\eta^{\star}=D/G_{\max}$ balances the two terms and yields
\[
\frac{1}{T}\sum_{k=0}^{T-1}\bigl(f(w_k)-f(w^{\star})\bigr)
\le
\frac{DG_{\max}}{\sqrt{T}} .
\]
Thus AdaGrad attains the optimal $O(1/\sqrt{T})$ rate for convex (smooth) problems with an explicit constant $DG_{\max}$.

\section*{Special Cases}
\begin{itemize}
\item \textbf{Sparse gradients.} When only a subset of coordinates is non‑zero at each iteration, each coordinate accumulates far fewer squared gradients, making the corresponding denominator grow slowly. The bound improves to $O\!\bigl(D\sqrt{\sum_{j=1}^{d}\log(1+G_{T}^{(j)})}/\sqrt{T}\bigr)$, matching the classic per‑coordinate AdaGrad analysis.
\item \textbf{Deterministic gradients.} Setting $g_k=\nabla f(w_k)$ removes all expectations; the same deterministic bound (3) holds.
\item \textbf{Strongly convex $f$.} If $f$ is $\mu$‑strongly convex, a refined analysis (using the strong‑convexity inequality) yields an $O\bigl((\log T)/T\bigr)$ rate, analogous to standard AdaGrad results for strongly convex objectives.
\end{itemize}

\end{document}

% ===== CORRECTION SUMMARY =====
% 1. [Step 1] Issue: index mismatch between algorithm (√(G_k+ε)) and proof (√(G_{k+1}+ε)). Fixed by using √(G_k+ε) consistently.
% 2. [Step 7] Issue: Lemma 2 was applied to the wrong denominator. Replaced with a correct inequality (13) using the small‑step‑size condition η≤√ε, and retained Lemma 2 for the later sum over 1/√(G_k+ε).
% 3. [Step 8] Issue: unjustified replacement of √(G_{k+1}+ε) by √(G_k+ε). Fixed by the argument in (13) and keeping the same index throughout.
% 4. [Step 9] Issue: division by zero at k=0. Handled by separating the k=0 term and starting the auxiliary sum at k=1; the bound absorbs the first term into constants.
% 5. Stochastic case: inserted expectation operators (Step 5‑6) and used Assumption 4 (unbiasedness). Smoothness (Assumption 2) is noted but not required for the convex rate.